\section{Макромир}
\label{sec:macro-known}

Макромир есть эти же уравнения, записанные для непрерывных полей.

Из \eqref{eq:second-law} и \eqref{eq:third-law} строится термодинамика.
Температура есть производная энергии по энтропии,
\begin{equation}
  \tempT = \Bigl(\frac{\partial \energyE}{\partial \entropyS}\Bigr)_{V,N},
  \label{eq:temperature}
\end{equation}
и из неё же --- потенциалы.
Интервал \eqref{eq:interval} один и тот же во всех инерциальных системах,
поэтому энергия и импульс составляют четырёхвектор
\begin{equation}
  p^{\mu} = (\energyE/\speedC,\,\vecP).
  \label{eq:four-momentum}
\end{equation}

Из этих уравнений собрана система единиц.
Числа~$\speedC$, $\planckHbar$, $\constBoltzmann$, заряд~$\chargeE$, массы частиц и постоянная тонкой структуры~$\alphaFs$ в ней измерены.
Сами уравнения их не выводят.
На этом этаже такое ещё терпимо: число может быть константой прибора.
Один этаж ниже тот же закон, если он претендует быть фундаментальным, уже не имеет права оставлять число вписанным рукой.

